\section{Tło projektu}

\subsection{Przegląd podobnych rozwiązań}

Detekcja tekstu generowanego przez modele językowe (LLM) jest aktywnie rozwijanym obszarem. Istnieją zarówno narzędzia komercyjne, jak i metody badawcze:

\begin{itemize}
	\item \textbf{GPTZero, Turnitin, Originality.ai} - systemy komercyjne stosowane m.in.\ w edukacji; opierają się na cechach stylometrycznych i modelach klasyfikacji.
	\item \textbf{DetectGPT / Fast-DetectGPT} - metody oparte na zaburzeniach prawdopodobieństwa (perturbation) tekstu względem modelu językowego.
	\item \textbf{Klasyfikatory nadzorowane} - modele typu RoBERTa/BERT trenowane na parach tekstów ludzkich i sztucznych (np.\ HC3, M4).
	\item \textbf{Watermarking} - oznaczanie wyjścia LLM na etapie generacji (np.\ schemat Kirchenbauer i in.); wymaga współpracy po stronie generatora.
\end{itemize}

W praktyce dokładność detektorów spada przy krótkich tekstach, parafrazie oraz tekstach wielojęzycznych. Dla języka polskiego liczba gotowych narzędzi i zbiorów jest wyraźnie mniejsza niż dla angielskiego.

\subsection{Stan wiedzy}

Problem sprowadza się zwykle do klasyfikacji binarnej: tekst ludzki vs tekst wygenerowany. Stosowane są dwa główne nurty:

\begin{enumerate}
	\item \textbf{Podejścia statystyczne / stylometryczne} - entropia, perplexity, burstiness, rozkład długości zdań, częstości n-gramów. Działają bez dużego zbioru treningowego, ale bywają wrażliwe na domenę i długość tekstu.
	\item \textbf{Uczenie maszynowe / głębokie} - klasyfikatory na embeddingach lub fine-tuning transformerów. Wymagają danych treningowych; przy małej liczbie przykładów łatwo o przeuczenie.
\end{enumerate}

W literaturze podkreśla się, że metryki oparte na perplexity zawodzą przy bardzo krótkich próbkach (kilka zdań). Różnorodność generatorów (GPT, Llama, Gemini itd.) utrudnia generalizację jednego modelu detekcji.

\subsection{Uwarunkowania komercyjne}

Rynek detekcji AI jest związany głównie z edukacją, wydawnictwami i moderacją treści. Dostawcy komercyjni rzadko publikują pełne wyniki i ograniczenia modeli. Dla projektu akademickiego istotne są:

\begin{itemize}
	\item dostępność otwartych zbiorów danych (Kaggle, Hugging Face),
	\item koszty API generatorów przy tworzeniu własnych próbek,
	\item brak gwarancji, że wynik detektora ma moc dowodową (ryzyko fałszywych oskarżeń).
\end{itemize}

\subsection{Wyniki konsultacji z ekspertem}

Konsultacje przeprowadzono z \textbf{dr.\ hab.\ Piotrem Hońko} (Wydział Informatyki PB). Poniżej pytania i odpowiedzi eksperta.

\begin{enumerate}
	\item \textit{Czy lepiej używać podejść statystycznych czy uczenia maszynowego?}\\
	\textbf{Odpowiedź:} Nie ma lepszego rozwiązania; dla małych zestawów danych podejście statystyczne może być lepsze.

	\item \textit{Czy teksty muszą być podobne tematycznie i czy w ogóle to jest możliwe?}\\
	\textbf{Odpowiedź:} Dla statystyki nie; dla ML tak.

	\item \textit{Jak podejść do problemu długości tekstu? Wiele metryk zawodzi na bardzo krótkich tekstach (np.\ 2-3 zdania).}\\
	\textbf{Odpowiedź:} Dobrać odpowiednie podejście.

	\item \textit{Czy warto nakładać minimalny limit długości tekstu w projekcie?}\\
	\textbf{Odpowiedź:} To zależy od podejścia - statystycznego albo ML.
\end{enumerate}

\subsection{Wyniki konsultacji AI}

Te same pytania zadano modelowi językowemu (ChatGPT). Odpowiedzi AI:

\begin{enumerate}
	\item \textit{Podejścia statystyczne vs ML:}\\
	\textbf{Odpowiedź AI:} Przy ograniczonej liczbie przykładów warto zacząć od cech statystycznych (perplexity, burstiness, stylometria). ML/transformers mają potencjał wyższej dokładności, ale potrzebują większego, zróżnicowanego zbioru. Sensowna jest ścieżka hybrydowa: najpierw baseline statystyczny, potem klasyfikator ML.

	\item \textit{Podobieństwo tematyczne tekstów:}\\
	\textbf{Odpowiedź AI:} Homogeniczność tematyczna pomaga modelom ML (mniej zmiennych zakłócających). Metody statystyczne są bardziej odporne na zmianę domeny, ale i tak warto kontrolować źródło tekstów. Pełna niezależność od tematu jest trudna do osiągnięcia.

	\item \textit{Problem krótkich tekstów:}\\
	\textbf{Odpowiedź AI:} Dla 2-3 zdań należy unikać metryk wymagających stabilnego oszacowania rozkładu tokenów. Lepiej agregować cechy na poziomie dokumentu o ustalonej minimalnej długości albo stosować modele specjalnie do krótkich sekwencji. Alternatywnie - klasyfikować dłuższe fragmenty (akapity).

	\item \textit{Minimalny limit długości:}\\
	\textbf{Odpowiedź AI:} Tak, warto wprowadzić próg (np.\ liczba słów/tokenów), zwłaszcza przy metrykach statystycznych. Dla ML próg też pomaga, ale model może działać na krótszych próbkach kosztem pewności. Limit powinien wynikać z eksperymentów na wybranym zbiorze.
\end{enumerate}

\subsection{Porównanie odpowiedzi eksperta i AI}

\begin{table}[H]
\centering
\small
\begin{tabularx}{\textwidth}{|p{2.6cm}|X|X|}
\hline
\textbf{Pytanie} & \textbf{dr hab.\ Piotr Hońko} & \textbf{Sztuczna inteligencja} \\
\hline
Statystyka vs ML &
Brak jednego lepszego rozwiązania; przy małych danych statystyka może być lepsza. &
Zacząć od baseline'u statystycznego; ML przy większych danych; możliwa hybryda. \\
\hline
Podobieństwo tematyczne &
Dla statystyki nie jest konieczne; dla ML tak. &
Homogeniczność pomaga ML; statystyka odporniejsza na domenę; pełna niezależność trudna. \\
\hline
Krótkie teksty &
Dobrać odpowiednie podejście do problemu. &
Unikać metryk niestabilnych na krótkich próbkach; agregować dłuższe fragmenty. \\
\hline
Limit długości &
Zależy od wybranego podejścia (statystyka / ML). &
Warto wprowadzić próg; dla statystyki szczególnie; próg ustalić eksperymentalnie. \\
\hline
\end{tabularx}
\caption{Porównanie konsultacji eksperckiej i AI}
\end{table}

Wnioski z konsultacji: przy niepewności co do wielkości zbioru i realnej wykonalności zadania rozsądne jest rozpoczęcie od metod prostszych (statystycznych), z możliwością przejścia do ML, oraz świadome ograniczenie zakresu (m.in.\ długość tekstu).
